\section{Pembahasan}
\subsection{Sejarah dan Evolusi Perkembangan Komputer}
\subsubsection{Pengertian Komputer}
Menurut \textcite{AYUBWIMATRA2018}, komputer adalah perangkat elektronik yang bekerja melalui empat tahap, yaitu menerima data sebagai masukan (\textit{input}), mengolahnya secara cepat dan akurat berdasarkan instruksi program yang tersimpan di dalam memori (\textit{processing}), menyimpan data serta hasil pengolahannya (\textit{storage}), dan menyajikan hasilnya berupa informasi (\textit{output}).

Hamacher, Vranesic, dan Zaky mendefinisikan komputer sebagai alat hitung elektronik berkecepatan tinggi yang menerima informasi digital sebagai masukan, memprosesnya sesuai program yang tersimpan di dalam memori (\textit{stored program}), dan menghasilkan keluaran berupa informasi \parencite[dikutip dalam][]{Yahfizham2019}.

Agar komputer dapat berfungsi untuk mengolah data menjadi informasi, harus ada tiga elemen utama yang saling terintegrasi:

\begin{enumerate}
    \item \textit{Hardware} (Perangkat Keras): Komponen fisik yang dapat dilihat dan diraba secara langsung, seperti CPU, RAM, \textit{motherboard}, \textit{keyboard}, dan monitor.
    \item \textit{Software} (Perangkat Lunak): Kumpulan program, instruksi, atau data digital yang mengatur kerja \textit{hardware} agar dapat menjalankan perintah tertentu.
    \item \textit{Brainware} (Pengguna/Manusia): Manusia yang mengoperasikan, mengatur, dan mengendalikan sistem komputer.
\end{enumerate}

\subsubsection{Alat Hitung Tradisional}

\subsubsection{Generasi Pertama: Tabung Vakum}


\subsubsection{Generasi Kedua: Transistor}
\subsubsection{Generasi Ketiga: \textit{Integrated Circuit (IC)}}
\subsubsection{Generasi Keempat: \textit{Microprocessor}}
\subsubsection{Generasi Kelima dan Perkembangan Terkini}
\subsubsection{Perbandingan antargenerasi}

\subsection{Komponen Utama Sistem Komputer dan Fungsinya}
\subsubsection{Gambaran Umum: \textit{hardware},\textit{Software},dan \textit{Brainware}}
\subsubsection{\textit{Central Processing Unit}: Unit Aritmatika dan Logika, Alat Kendali, dan \textit{Register}}
\subsubsection{Memori: Memori Utama (RAM, ROM) dan Penyimpanan Sekunder}
\subsubsection{Perangkat Masukan dan Keluaran}

\subsection{Hubungan Antarkomponen Sistem Komputer}
\subsubsection{Arsitektur Von Neumann}
\subsubsection{\textit{Bus} Sistem}
\subsubsection{Alur Kerja Antarkomponen}

\subsection{Pengaruh Perkembangan Teknologi terhadap Komponen}
\subsubsection{Perubahan pada CPU}
\subsubsection{Perubahan pada Memori dan Penyimpanan}
\subsubsection{Perubahan pada Perangkat Input dan Output}
\subsubsection{Dampak Keseluruhan}